\section{Diseño del Algoritmo: Relajación Bidireccional y Poda}

Para romper la dependencia circular sin alterar la optimalidad global, analizamos la estructura geométrica de cualquier camino óptimo dentro de una misma capa $i$.

\subsection{Descomposición Bidireccional por Monotonía Horizontal}

\textbf{Lema 1 (Monotonía Horizontal entre capas):} \textit{En un digrafo triangular estratificado con pesos positivos $v_{i,j} > 0$, ningún camino óptimo puede cambiar de dirección horizontal dentro de una misma fila $i$.}

\textbf{Demostración (Por reducción al absurdo):} \\
Supongamos que existe un camino óptimo $P^*$ que, encontrándose en la fila $i$, realiza un movimiento hacia la derecha $(i, j) \to (i, j+1)$ e inmediatamente (o tras varios pasos horizontales) invierte su sentido hacia la izquierda regresando a $(i, j)$. Dicho subcamino tendría la forma $\langle \dots, (i, j), \dots, (i, j+k), \dots, (i, j), \dots \rangle$ con $k \geq 1$. Dado que $v_{i,m} > 0$ para todo vértice, si recortamos el ciclo intermedio reemplazando todo el bucle por una única visita al nodo $(i, j)$, obtenemos un camino válido $P'$ cuyo costo verifica:
\[
    C(P') = C(P^*) - \sum_{\text{ciclo}} v_{i,m} < C(P^*)
\]
Esto contradice la hipótesis de que $P^*$ era óptimo. El razonamiento es idéntico si primero se desplaza hacia la izquierda y luego hacia la derecha. $\blacksquare$

\vspace{1ex}
\noindent\textbf{Corolario (Descomposición en Dos Fases Acíclicas):} \\
Todo segmento horizontal de un camino óptimo en la fila $i$ es estrictamente monótono: o bien ingresa a la fila $i$ por una columna $k \geq j$ y se desplaza exclusivamente hacia la izquierda ($\leftarrow$), o bien ingresa por una columna $k \leq j$ y se desplaza exclusivamente hacia la derecha ($\rightarrow$). Esto permite descomponer el cálculo de la fila $i$ en dos subgrafos acíclicos independientes:
\begin{enumerate}
    \item \textbf{Estado Auxiliar Estrella ($\mathcal{OPT}^*[i][j]$):} Representa el costo mínimo para llegar a $(i,j)$ permitiendo únicamente descensos desde la fila $i-1$ y movimientos horizontales \textit{hacia la izquierda}. Se calcula yendo de la fila de derecha a izquierda.
    \item \textbf{Estado Definitivo ($\mathcal{OPT}[i][j]$):} Consolida el costo mínimo absoluto comparando el mejor camino proveniente de la derecha ($\mathcal{OPT}^*[i][j]$) contra la opción de venir desde el vecino izquierdo ya consolidado ($\mathcal{OPT}[i][j-1] + v_{i,j}$). Se calcula yendo de izquierda a derecha.
\end{enumerate}

\subsection{Teorema de Poda por Dominancia Geométrica en las Fronteras}

A diferencia de una grilla rectangular donde el doble barrido debe evaluar los $m$ elementos completos en ambas direcciones, demostramos que la geometría triangular induce una relación de dominancia estricta en ambas esquinas de cada fila $i \geq 2$.

\textbf{Teorema 1 (Poda de Esquinas y Reducción de Estados Estrella):} \textit{Para toda capa $i \geq 2$, los vértices extremos $p = 0$ (esquina izquierda) y $u = i$ (esquina derecha) nunca pueden formar parte de un camino óptimo que se desplace horizontalmente hacia el interior de la fila $i$. En consecuencia: (a) el pre-cálculo de $\mathcal{OPT}^*$ en el anteúltimo elemento $a = i-1$ prescinde de evaluar a $u = i$, y (b) el segundo elemento $d = 1$ verifica $\mathcal{OPT}[i][1] = \mathcal{OPT}^*[i][1]$ sin evaluar a $p = 0$.}

\textbf{Demostración Algebraica:}
\begin{itemize}
    \item \textbf{Parte (a) — Dominancia en la Frontera Derecha ($u = i$ y $a = i-1$):} \\
    Por la topología triangular, el extremo derecho $(i, i)$ carece de padre superior derecho, la fila $i-1$ solo llega hasta la columna $i-1$. Su único padre en la capa superior es $(i-1, i-1)$. Si un camino desciende desde $(i-1, i-1)$ hacia la esquina $(i, i)$ y luego retrocede horizontalmente hacia el anteúltimo nodo $(i, i-1)$, su costo acumulado es:
    \[
        C_{\text{esquina}} = \mathcal{OPT}[i-1][i-1] + v_{i,i} + v_{i,i-1}
    \]
    Sin embargo, el vértice $(i-1, i-1)$ es simultáneamente padre directo de $(i, i-1)$ mediante una arista descendente. El costo de bajar directamente de $(i-1, i-1)$ a $(i, i-1)$ es:
    \[
        C_{\text{directo}} = \mathcal{OPT}[i-1][i-1] + v_{i,i-1}
    \]
    Como $v_{i,i} > 0$, se cumple invariablemente que $C_{\text{directo}} < C_{\text{esquina}}$. Por lo tanto, la transición $(i, i) \to (i, i-1)$ es estrictamente subóptima y se poda del cálculo, definiendo el caso base derecho de $\mathcal{OPT}^*$ como:
    \[
        \mathcal{OPT}^*[i][a] = v_{i,a} + \min\big(\mathcal{OPT}[i-1][a-1], \; \mathcal{OPT}[i-1][a]\big)
    \]

    \item \textbf{Parte (b) — Dominancia en la Frontera Izquierda ($p = 0$ y $d = 1$):} \\
    De manera simétrica, el extremo izquierdo $(i, 0)$ tiene un único padre superior: $(i-1, 0)$. Cualquier camino que descienda a $(i, 0)$ y luego avance horizontalmente hacia la derecha al segundo nodo $(i, 1)$ da un costo $\mathcal{OPT}[i-1][0] + v_{i,0} + v_{i,1}$. Dado que $(i-1, 0)$ también conecta directamente con $(i, 1)$ con costo $\mathcal{OPT}[i-1][0] + v_{i,1}$ (el cual ya fue evaluado dentro de $\mathcal{OPT}^*[i][1]$), y siendo $v_{i,0} > 0$, el paso desde $(i,0)$ hacia $(i,1)$ queda estrictamente dominado. En consecuencia, el segundo elemento consolida su óptimo directamente sin comparaciones adicionales:
    \[
        \mathcal{OPT}[i][1] = \mathcal{OPT}^*[i][1]
    \]
\end{itemize}
Gracias a las dos podas, para una fila de $k = i+1$ elementos solo se requieren calcular $k - 2 = i - 1$ estados estrella (desde $j = i-1$ hasta $j = 1$). $\blacksquare$

\vspace{1ex}
\noindent\textbf{Resumen de la Ecuación de Recurrencia Final (Para toda capa $i \geq 2$):} \\
Definiendo los punteros notables $p = 0$, $d = 1$, $u = i$ y $a = i-1$:
\[
    \text{\textbf{Fase 1 ($\mathcal{OPT}^*$):}} \quad
    \begin{cases}
        \mathcal{OPT}^*[i][a] = v_{i,a} + \min\big(\mathcal{OPT}[i-1][a-1], \; \mathcal{OPT}[i-1][a]\big) \\
        \mathcal{OPT}^*[i][j] = v_{i,j} + \min\big(\mathcal{OPT}[i-1][j-1], \; \mathcal{OPT}[i-1][j], \; \mathcal{OPT}^*[i][j+1]\big) & \forall j \in [a-1 \dots 1]
    \end{cases}
\]
\[
    \text{\textbf{Fase 2 ($\mathcal{OPT}$):}} \quad
    \begin{cases}
        \mathcal{OPT}[i][p] = v_{i,p} + \min\big(\mathcal{OPT}[i-1][p], \; \mathcal{OPT}^*[i][d]\big) \\
        \mathcal{OPT}[i][d] = \mathcal{OPT}^*[i][d] \\
        \mathcal{OPT}[i][j] = \min\big(v_{i,j} + \mathcal{OPT}[i][j-1], \; \mathcal{OPT}^*[i][j]\big) & \forall j \in [2 \dots u-1] \\
        \mathcal{OPT}[i][u] = v_{i,u} + \min\big(\mathcal{OPT}[i-1][u-1], \; \mathcal{OPT}[i][u-1]\big)
    \end{cases}
\]

\subsection{Optimización Espacial In-Place e Implementación}\label{rec:codigo}

En una primera aproximación podrían instanciarse dos matrices cuadradas de $n \times n$ (\verb|ESTRELLA| y \verb|OPT|), consumiendo $2n^2$ celdas en memoria. Sin embargo, observamos dos propiedades clave que permiten reducir drásticamente el consumo espacial:
\begin{enumerate}
    \item \textbf{Dimensionamiento Escalonado:} La fila $i$ solo posee $i+1$ columnas válidas. Al reservar listas de longitud exacta $i+1$, reducimos el espacio base de $n^2$ a $\frac{n(n+1)}{2} \approx 0.5n^2$.
    \item \textbf{Sobrescritura In-Place (Eliminación total de \texttt{ESTRELLA}):} Los valores de $\mathcal{OPT}^*[i]$ solo se utilizan durante el cálculo de la misma fila $i$. Por lo tanto, en la Fase 1 escribimos los valores de $\mathcal{OPT}^*$ directamente dentro del arreglo definitivo \verb|OPT[i]| de derecha a izquierda. Después, en la Fase 2, recorremos \verb|OPT[i]| de izquierda a derecha actualizando \verb|OPT[i][j] = min(OPT[i][j], piramide[i][j] + OPT[i][j-1])|, donde el valor previo que estaba en \verb|OPT[i][j]| es exactamente $\mathcal{OPT}^*[i][j]$. Esto elimina el $100\%$ de la matriz auxiliar \verb|ESTRELLA|, logrando un ahorro total del $75\%$ de memoria respecto al planteo matricial ingenuo.
\end{enumerate}

A continuación se presenta la implementación completa en lenguaje Python:

\begin{lstlisting}[language=Python, caption={Implementación del algoritmo de PD Bidireccional In-Place y su reconstrucción.}]
def camino_minimo_piramide(piramide: list[list[int]]) -> tuple[int, list[list[int]]]:
    if not piramide:
        return 0, []
    n = len(piramide)
    OPT = [[0] * (i + 1) for i in range(n)]  # Matriz escalonada unica (sin tabla ESTRELLA)

    # Casos base (filas 0 y 1)
    OPT[0][0] = piramide[0][0]
    if n == 1:
        return OPT[0][0], [[0, 0]]
    OPT[1][0] = piramide[1][0] + OPT[0][0]
    OPT[1][1] = piramide[1][1] + OPT[0][0]

    for i in range(2, n):
        p, d, u, a = 0, 1, i, i - 1

        # Fase 1: Barrido Derecha a Izquierda (OPT* almacenado in-place sobre OPT[i])
        OPT[i][a] = piramide[i][a] + min(OPT[i - 1][a - 1], OPT[i - 1][a])
        for j in range(a - 1, 0, -1):
            OPT[i][j] = piramide[i][j] + min(OPT[i - 1][j - 1], OPT[i - 1][j], OPT[i][j + 1])

        # Fase 2: Barrido Izquierda a Derecha (Consolidacion de OPT definitivo)
        OPT[i][p] = piramide[i][p] + min(OPT[i - 1][p], OPT[i][d])
        # OPT[i][d] ya contiene OPT*[i][d], queda consolidado por Teorema 1(b)
        for j in range(2, u):
            OPT[i][j] = min(OPT[i][j], piramide[i][j] + OPT[i][j - 1])
        OPT[i][u] = piramide[i][u] + min(OPT[i - 1][u - 1], OPT[i][u - 1])

    ult = n - 1
    min_valor = min(OPT[ult])
    min_col = OPT[ult].index(min_valor)
    return min_valor, _reconstruccion(piramide, OPT, [min_valor, ult, min_col])

def _reconstruccion(piramide: list[list[int]], OPT: list[list[int]], destino: list) -> list[list[int]]:
    if destino[0] == float('inf'):
        return []
    camino = []
    fila, col = destino[1], destino[2]
    camino.append([fila, col])

    while fila > 0:
        buscado = OPT[fila][col] - piramide[fila][col]
        if col > 0 and OPT[fila - 1][col - 1] == buscado:
            fila, col = fila - 1, col - 1
        elif col < fila and OPT[fila - 1][col] == buscado:
            fila -= 1
        elif col > 0 and OPT[fila][col - 1] == buscado:
            col -= 1
        elif col < fila and OPT[fila][col + 1] == buscado:
            col += 1
        camino.append([fila, col])

    return camino[::-1]
\end{lstlisting}

